% Generated draft: requires mathematical and source review.
\documentclass[11pt]{article}
\usepackage[margin=25mm]{geometry}
\usepackage{fontspec}
\setmainfont{Latin Modern Roman}
\usepackage{amsmath,amssymb,mathtools,unicode-math}
\setmathfont{Latin Modern Math}
\usepackage{enumitem,needspace,xurl,hyperref}
\setlength{\parindent}{0pt}
\setlength{\parskip}{5pt}
\setlength{\emergencystretch}{3em}
\newcommand{\N}{\mathbb{N}}
\newcommand{\Z}{\mathbb{Z}}
\newcommand{\Q}{\mathbb{Q}}
\newcommand{\R}{\mathbb{R}}
\newcommand{\C}{\mathbb{C}}
\newcommand{\F}{\mathbb{F}}
\newcommand{\A}{\mathbb{A}}
\newcommand{\PP}{\mathbb{P}}
\newcommand{\E}{\mathbb{E}}
\newcommand{\eps}{\varepsilon}
\newcommand{\dd}{\,\mathrm{d}}
\begin{document}
\section*{UnsolvedMath ID: 6700086}
\textbf{Source problem code:} catalogue entry
\subsection{Definitions and mathematical statement}
The extracted ``Hyperbolic Volume Inequality'' begins with ``Then,'' but its antecedent hypotheses on the Riemannian manifolds $X$ and $X_0$ are missing. Write $n=\dim X$ where appropriate. Relative to those missing hypotheses, the proposed conclusion is that every continuous $f_0:X\to X_0$ is homotopic to a map $f$ satisfying
\[
 \operatorname{vol}_n(f(X))\le\operatorname{vol}(X).
\]
Moreover the map can be chosen with strict inequality unless both $X$ has constant negative curvature and $f_0$ is homotopic to a locally isometric map. Thus outside that exceptional class the stated target is
\[
 \exists f\simeq f_0,\qquad
 \operatorname{vol}_n(f(X))<\operatorname{vol}(X).
\]
The ``unless'' clause does not assert that every exceptional map attains equality.

The volume of the image could mean its Hausdorff/Riemannian measure as a subset, or a parametrized volume counting multiplicity. Those are different for maps of nontrivial degree, and the source does not choose. The dimensions, compactness, metrics and curvature assumptions, finiteness of volume, and regularity of the representative $f$ are likewise absent. No scalar-curvature lower bound or target-curvature hypothesis is inferred merely from the title.

The passage cites the constant-sectional-curvature-$-1$ case for both manifolds as known via simplicial volume, describing it as a generalization of Mostow rigidity. For closed oriented $n$-manifolds, the usual simplicial volume $\|X\|$ is the infimum of $\sum_j|a_j|$ over real singular cycles $\sum_ja_j\sigma_j$ representing $[X]$. The source recalls degree monotonicity and covering equality. With nonnegative degree $d$ its displayed form is
\[
 \|X\|\ge d\,\|Y\|,
 \qquad \|X\|=d\,\|Y\|
       \text{ for a }d\text{-sheeted covering}.
\]
For signed degrees the usual monotonicity uses $|d|$; the source does not specify a sign convention. It also reports nonzero simplicial volume when $X$ admits negative sectional curvature. Its wording says compact orientable manifolds without a boundary convention; a relative version would be required in the presence of boundary. These contextual claims do not supply the missing hypotheses or volume convention for the main conjecture.
\subsection{Short English statement}
Find a homotopic representative with image volume at most the domain volume, with strict improvement outside the stated constant-negative-curvature/local-isometry exception. The intended ambient hypotheses are lost from the extraction, so the formula cannot be asserted for arbitrary manifolds.

The supplied summary reports the closed hyperbolic special case via Gromov--Thurston simplicial volume and calls the general formulation open. It supplies no precise general theorem or missing antecedent. The partial-progress label and 2017 passage are retained with current status marked for review.
\subsection{Sources}
\begin{itemize}
\item Local UnsolvedMath registry, permanent ID 6700086, source code catalogue entry, scalar-curvature Question [?90], ``Hyperbolic Volume Inequality,'' proposed by Misha Gromov, 2017. Dataset curation and metadata: Ulam AI and UnsolvedMath Contributors, CC BY 4.0.
\item Misha Gromov, \emph{101 Questions, Problems and Conjectures around Scalar Curvature} (2017), printed tag [?90], occurrence 89, PDF page 88, \url{https://www.ihes.fr/~gromov/wp-content/uploads/2018/08/101-problemsOct1-2017.pdf}; normalized pdf-text extraction, accessed 2026-07-29.
\item Mostow rigidity and simplicial volume are named in the source; the summary also names Gromov--Thurston. No complete bibliographic entries or exact theorem hypotheses are supplied.
\item Source-specific provenance: the author-hosted PDF is incomplete and publicly available; current status and statement-reuse terms NEEDS\_REVIEW, separately from the dataset licence.
\end{itemize}
\end{document}
